\chapter{Plate Equation}

\section*{Introduction}

The plate equation governs the bending of thin elastic plates under transverse loads. Formulated by Augustus Edward Hough Love in 1888 as part of the Kirchhoff--Love theory of plates, it is a fourth-order partial differential equation that generalizes the one-dimensional beam-bending equation (Euler--Bernoulli) to two dimensions. The equation is fundamental to structural engineering, geophysics (lithospheric flexure), and the design of pressure vessels, floors, and microelectromechanical systems (MEMS).

A thin plate is a flat structural element whose thickness is much smaller than its lateral dimensions. When a transverse load $q$ (pressure per unit area) is applied, the plate deflects by an amount $w(x,y)$. The restoring force comes from the plate's bending stiffness $D = Eh^{3}/[12(1-\nu^{2})]$, where $E$ is Young's modulus, $h$ is the thickness, and $\nu$ is Poisson's ratio. The biharmonic operator $\nabla^{4}w$ relates the fourth spatial derivatives of the deflection to the applied load, encoding the resistance of the plate to bending in two dimensions simultaneously.


\begin{equation}
\boxed{\; D\nabla^{4}w = q\;}
\label{eq:plate}
\end{equation}

where $D = Eh^{3}/[12(1-\nu^{2})]$ is the flexural rigidity and $\nabla^{4} = \nabla^{2}\nabla^{2}$ is the biharmonic operator.

\section*{Fundamental Equations Used}

The derivation starts from the definition of bending moments $M_{ij} = \int \sigma_{ij}z\,dz$ as integrals of stress through the plate thickness.

\section*{Assumptions}

\textbf{Assumptions:} The plate is thin ($h \ll a$, where $a$ is the lateral dimension); the deflection is small ($w \ll h$); the plate is homogeneous, isotropic, and linearly elastic; transverse shear deformation is negligible.


\textbf{Limitations:} Does not account for large deflections (von Kármán plate theory is needed for $w \sim h$); does not include shear deformation (Mindlin--Reissner theory for thicker plates); not valid for composite or laminated plates without modification.


\section*{Mathematical Derivation}

\textbf{Step 1: Define the plate geometry and bending moments.}

Consider a thin plate of thickness $h$ in the $xy$-plane. The deflection $w(x,y)$ is small compared to $h$. The bending moments per unit length are defined as:
\begin{equation}
M_{xx} = \int_{-h/2}^{h/2} \sigma_{xx}\,z\,dz, \qquad M_{yy} = \int_{-h/2}^{h/2} \sigma_{yy}\,z\,dz, \qquad M_{xy} = \int_{-h/2}^{h/2} \sigma_{xy}\,z\,dz,
\end{equation}
where $\sigma_{ij}$ are the in-plane stresses and $z$ is measured from the midplane.

\textbf{Step 2: Relate stresses to curvatures using the constitutive law.}

For a linearly elastic, isotropic plate, the strain--dispatibility relations (Kirchhoff hypothesis) give the in-plane strains as:
\begin{equation}
\varepsilon_{xx} = -z\,\frac{\partial^{2}w}{\partial x^{2}}, \qquad \varepsilon_{yy} = -z\,\frac{\partial^{2}w}{\partial y^{2}}, \qquad \varepsilon_{xy} = -z\,\frac{\partial^{2}w}{\partial x\,\partial y}.
\end{equation}
Using Hooke's law for plane stress ($\sigma_{zz} \approx 0$):
\begin{equation}
\sigma_{xx} = \frac{E}{1-\nu^{2}}\bigl(\varepsilon_{xx} + \nu\varepsilon_{yy}\bigr), \qquad \sigma_{yy} = \frac{E}{1-\nu^{2}}\bigl(\varepsilon_{yy} + \nu\varepsilon_{xx}\bigr).
\end{equation}

\textbf{Step 3: Integrate to find the bending moments.}

Substituting the strains into the moment integrals:
\begin{equation}
M_{xx} = -\frac{E}{1-\nu^{2}}\!\left(\frac{\partial^{2}w}{\partial x^{2}} + \nu\,\frac{\partial^{2}w}{\partial y^{2}}\right)\!\int_{-h/2}^{h/2} z^{2}\,dz = -D\!\left(\frac{\partial^{2}w}{\partial x^{2}} + \nu\,\frac{\partial^{2}w}{\partial y^{2}}\right),
\end{equation}
where $D = Eh^{3}/[12(1-\nu^{2})]$ is the flexural rigidity. Similarly:
\begin{equation}
M_{yy} = -D\!\left(\frac{\partial^{2}w}{\partial y^{2}} + \nu\,\frac{\partial^{2}w}{\partial x^{2}}\right), \qquad M_{xy} = -D(1-\nu)\,\frac{\partial^{2}w}{\partial x\,\partial y}.
\end{equation}

\textbf{Step 4: Apply force and moment equilibrium to a plate element.}

Consider a small element $dx \times dy$. Vertical force equilibrium gives:
\begin{equation}
\frac{\partial Q_x}{\partial x} + \frac{\partial Q_y}{\partial y} = -q,
\end{equation}
where $Q_x$ and $Q_y$ are shear forces per unit length and $q$ is the transverse load per unit area. Moment equilibrium (summing moments about the $y$-axis and $x$-axis) gives:
\begin{equation}
Q_x = \frac{\partial M_{xx}}{\partial x} + \frac{\partial M_{xy}}{\partial y}, \qquad Q_y = \frac{\partial M_{yy}}{\partial y} + \frac{\partial M_{xy}}{\partial x}.
\end{equation}

\textbf{Step 5: Combine to derive the biharmonic equation.}

Substituting the shear forces into the force equilibrium equation:
\begin{equation}
\frac{\partial^{2}M_{xx}}{\partial x^{2}} + 2\,\frac{\partial^{2}M_{xy}}{\partial x\,\partial y} + \frac{\partial^{2}M_{yy}}{\partial y^{2}} = -q.
\end{equation}
Inserting the moment--curvature relations:
\begin{equation}
D\!\left(\frac{\partial^{4}w}{\partial x^{4}} + 2\,\frac{\partial^{4}w}{\partial x^{2}\partial y^{2}} + \frac{\partial^{4}w}{\partial y^{4}}\right) = q.
\end{equation}
Recognizing the biharmonic operator $\nabla^{4}w = \partial^{4}w/\partial x^{4} + 2\,\partial^{4}w/\partial x^{2}\partial y^{2} + \partial^{4}w/\partial y^{4}$:
\begin{equation}
\boxed{\;D\,\nabla^{4}w = q.\;}
\end{equation}

\textbf{Step 6: Interpret the result.}

The equation states that the applied load is balanced by the resistance of the plate to bending. The fourth-order nature arises because the load is related to the divergence of shear forces, which are divergences of bending moments, which are proportional to second derivatives of deflection---giving four spatial derivatives in total.

\section*{Characteristics of the Equation}

\begin{itemize}
  \item \textbf{Fourth-order PDE:} The biharmonic equation requires four boundary conditions per edge (compared to two for the second-order membrane equation).
  \item \textbf{Bending moments:} $M_{xx} = -D(\partial^{2}w/\partial x^{2} + \nu\,\partial^{2}w/\partial y^{2})$; $M_{yy} = -D(\partial^{2}w/\partial y^{2} + \nu\,\partial^{2}w/\partial x^{2})$; $M_{xy} = -D(1-\nu)\partial^{2}w/\partial x\,\partial y$.
  \item \textbf{Stiffness:} $D \propto h^{3}$, so doubling the thickness increases stiffness eightfold.
  \item \textbf{Poisson effect:} The factor $(1-\nu^{2})$ in $D$ accounts for the biaxial constraint: a plate is stiffer than a beam of the same cross-section because transverse contraction is constrained.
\end{itemize}
\section*{Solution Methods}

\begin{enumerate}
  \item \textbf{Separation of variables:} For rectangular plates with simply supported edges, assume $w = \sum W_{mn}\sin(m\pi x/a)\sin(n\pi y/b)$ and solve for the Fourier coefficients.
  \item \textbf{Circular plates:} Use polar coordinates $(r,\theta)$ and the biharmonic in polar form; for axisymmetric loading, solutions involve $r^{2}\ln r$, $r^{2}$, and constants.
  \item \textbf{Numerical methods:} Finite element methods (FEM) discretize the plate into elements and solve the system $[K]\{w\} = \{q\}$.
  \item \textbf{Energy methods:} The Rayleigh--Ritz method minimizes the total potential energy (bending energy minus load work) to find approximate deflections.
\end{enumerate}
\section*{Verifications}

Plate theory predictions are verified by strain gauge measurements, deflection measurements under known loads, and comparison with finite element solutions. The deflection of simply supported rectangular plates under uniform load agrees with series solutions to within experimental uncertainty. The natural frequencies of vibrating plates (Chladni patterns) confirm the eigenvalue structure of the biharmonic operator.
\section*{Applications}

\begin{itemize}
  \item \textbf{Civil engineering:} Floor slabs, bridge decks, and flat roof structures are designed using plate theory.
  \item \textbf{Pressure vessels:} Flat end plates of cylindrical pressure vessels are analyzed as clamped plates under uniform pressure.
  \item \textbf{Geophysics:} The flexure of the lithosphere under volcanic loads (e.g., Hawaiian Islands) is modeled as a plate on an elastic foundation.
  \item \textbf{MEMS:} Micro-mechanical membranes and diaphragms in sensors and actuators are analyzed using the plate equation.
\end{itemize}
\section*{What Richard Feynman Would Have Asked}

\subsection*{Plain-English Translation}
\textit{``What does the plate equation say, if I describe it to an engineer who has never studied elasticity theory?''}

When you stand on a floor, it bends---just a tiny bit. The plate equation tells you exactly how much it bends, given how thick it is, what material it is made of, and how heavy you are. The bending stiffness $D$ captures the material's resistance to bending (thicker and stiffer materials bend less), and the biharmonic operator $\nabla^{4}w$ captures how the bending spreads out in two dimensions---unlike a beam, which bends in only one direction, a plate distributes the load sideways as well. The fourth-order nature of the equation means that the bending behavior is more complex than simple stretching: it involves the curvature of the curvature.

\subsection*{Localized Mechanics}
\textit{``At the level of a small element of the plate, what forces and moments are in balance?''}

Consider a tiny square element $dx \times dy$ of the plate. The transverse load $q\,dx\,dy$ pushes it down. Resisting this are shear forces on the four edges and bending moments (couples) that twist the element. The bending moment $M_{xx}$ is the torque per unit length about the $y$-axis; it is proportional to the curvature $\partial^{2}w/\partial x^{2}$ (with a Poisson correction from the $y$-curvature). The equilibrium equation for the element is obtained by summing all vertical forces and all moments. The result is $D\nabla^{4}w = q$, which says that the applied load is balanced by the divergence of the shear forces, which in turn are the divergence of the bending moments, which are proportional to the second derivatives of $w$. This double differentiation is why the equation is fourth-order.

\subsection*{Testing Extreme Limits}
\textit{``What happens when the plate is very thin, very thick, very large, or loaded very heavily?''}

For a very thin plate ($h \to 0$), $D \to 0$ and the plate offers no resistance to bending---it behaves like a membrane (if tensioned) or collapses. For a very thick plate ($h/a \gtrsim 0.1$), shear deformation becomes important and the Kirchhoff--Love theory underestimates deflections; the Mindlin--Reissner theory, which includes shear, must be used. For a very large plate ($a \to \infty$) with a localized load, the deflection spreads out and the plate behaves like a half-space (Boussinesq problem). For very heavy loads (large $q$ or $w \sim h$), nonlinear effects dominate: the plate enters the membrane regime where stretching, not bending, carries the load, and the von Kármán equations must be used. In the extreme limit of a point load on an infinite plate, the deflection is logarithmic in $r$ (like the 2D Green's function of the biharmonic operator), diverging at the origin---a mathematical idealization that in practice is regularized by the finite size of the load application area.

\subsection*{Dimensionality and Geometry}
\textit{``How does the geometry of the plate---its shape, curvature, and boundary conditions---affect the solutions?''}

The shape of the plate determines the eigenfunctions of the biharmonic operator. For a rectangular plate with simply supported edges, the eigenfunctions are $\sin(m\pi x/a)\sin(n\pi y/b)$ and the eigenvalues are $(m^{2}\pi^{2}/a^{2} + n^{2}\pi^{2}/b^{2})^{2}$. For a circular plate, the eigenfunctions involve Bessel functions and angular harmonics $e^{in\theta}$. For a clamped edge, $w = 0$ and $\partial w/\partial n = 0$ (zero deflection and zero slope); for a free edge, $M_{nn} = 0$ and $V_{n} = 0$ (zero moment and zero shear). The eigenvalues determine the natural frequencies of vibration: $\omega_{mn} = \sqrt{D/\rho h}\,\lambda_{mn}$ where $\lambda_{mn}$ is the eigenvalue. Chladni patterns---the nodal lines of vibrating plates---are visual representations of these eigenfunctions. For curved shells (doubly curved), the plate equation is replaced by the Donnell--Mushtari shell equations, which include membrane-bending coupling.

\subsection*{Cross-Disciplinary Connection}
\textit{``Where does the biharmonic equation $\nabla^{4}w = q$ appear outside of structural engineering?''}

The biharmonic equation appears in: (1) Stokes flow (low-Reynolds-number fluid mechanics), where the stream function $\psi$ satisfies $\nabla^{4}\psi = 0$ for 2D incompressible flow; (2) elasticity theory, where the Airy stress function for plane problems satisfies $\nabla^{4}\phi = 0$ in the absence of body forces; (3) image processing, where biharmonic smoothing (thin-plate splines) minimizes bending energy to interpolate scattered data; (4) geophysics, where the flexure of tectonic plates under volcanic or sediment loads is modeled by $\nabla^{4}w + kw/D = q/D$ (with a Winkler foundation term $kw$); and (5) quantum field theory, where the biharmonic operator appears in the effective action for certain conformal field theories. The unifying structure is that the biharmonic equation is the simplest fourth-order elliptic PDE, and it governs any system where the dominant energy is proportional to the square of the curvature.

% ============================================================
% CHAPTER 48 — Laplace's Equation
% ============================================================
\section*{FAQ}

\begin{enumerate}
\item \textbf{Why is the plate equation fourth-order while the membrane equation is second-order?}
A membrane (like a soap film) resists only in-plane tension; a plate resists bending, which involves curvatures (second derivatives). The moment--curvature relation introduces second derivatives of curvature, leading to fourth derivatives of deflection.

\item \textbf{What is the difference between a plate and a membrane?}
A plate has bending stiffness ($D > 0$) and can support transverse loads without in-plane tension. A membrane has no bending stiffness ($D = 0$) and can only support loads through in-plane tension---like a drumhead.

\item \textbf{What happens for large deflections?}
When $w \sim h$, the midplane stretches and membrane stresses become significant. The linear plate equation is replaced by the nonlinear von Kármán equations, which couple in-plane and out-of-plane deformations.
\end{enumerate}
\section*{Important Bibliography}

\begin{enumerate}
\item Timoshenko, S.P. and Woinowsky-Krieger, S., \textit{Theory of Plates and Shells}, 2nd ed. (McGraw-Hill, 1959), Chapter 2.
\item Love, A.E.H., \textit{A Treatise on the Mathematical Theory of Elasticity}, 4th ed. (Dover, 1944), Chapter 23.
\item Ugural, A.C., \textit{Stresses in Plates and Shells}, 4th ed. (McGraw-Hill, 1999), Chapter 3.
\item Ventsel, E. and Krauthammer, T., \textit{Thin Plates and Shells} (Marcel Dekker, 2001), Chapter 2.
\end{enumerate}

